\scanpage{19}
\mspage{4}
\unsure{uncer|under} $\underline R^{+*}$, which plays here the role of the parameter space of \inserted{all} possible (non degenerate) ``radii'' at $s_i$. If we admit also radius zero, we accordingly get a parameter space $R_i^{\wedge}$, which may be viewed as a torsor of sorts \unsure{under} $\underline R^+$. Thus the set of radii for a given admissible system of discs $D_i$ around the points $s_i$ may be viewed as a point of the product-space
\[
r=(r_i)_{i\in I}\in R^{\wedge}=\prod_{i\in I}R_i^{\wedge}\ .
\]
My expectation is that an admissible set of discs $(D_i)$ is well determined by the knowledge of the corresponding set\inserted{$r$} of radii \deleted{$\unsure{r}$}, and moreover that a given set\inserted{$r$} of radii \deleted{$\unsure{r}$} corresponds to \deleted{\unsure{a}}\inserted{an admissible} system of \deleted{admissible} discs iff it satisfies a set of unequalities
\[
r_i<\rho_i\ ,
\]
where
\[
\rho=(\rho_i)_{i\in I}\in R=\prod_{i\in I}R_i
\]
is a \inserted{some} fixed system of radii, corresponding to a fixed system of choices of hermitian metrics in the tangent spaces $T_{s_i}$.

I now see that this ``expectation'' does'nt quite match\inserted{with} the previous one, about a ``natural operation'' of $(\underline R^{\unsure{+*|+}})^I$ upon $\widetilde{MB}$, having certain properties - it would match only if all $\rho_i$ where equal to $+\infty$ (hence not in $R_i$ itself strictly speaking). I must confess I did'nt look too thoroughly yet at the situation, and moreover I've been busy with rather different kind of things for the last two or three months, and lost\inserted{contact} a little\ldots\ \deleted{contact.} What is clear however is that the main key to an understanding of the general situation, is in an understanding of the basic particular case of Thurston's pants. If we number $0,1,\infty$ the three ``holes'' of such a pant, the surface $X^{\wedge}$ can be identified canonically to the Riemann sphere, and the basic question then is to understand how the pant is imbedded in this sphere $\Sigma$, as a complement of \deleted{(open)} the union of (open) discs around the points $0,1,\infty$, these discs forming an ``admissible system''. So the main question is about understanding the structure of all possible admissible systems of three discs on $\Sigma$.

\unsure{Pyzzling|Puzzling} a little about this problem, the following model came to my mind (corresponding to ``limiting radii'' $\rho_i$ which are \uline{finite}, not infinite). I view $\Sigma$ as endowed with it's usual euclidean metric, for which the real projective line is a great circle, with $0,1,\infty$ at equal distance from each other on this equator. These points may be viewed as the centers of three ``orange \deleted{quarters}\inserted{slices}'', making up a cellular subdivision of $\Sigma$, where the common boundary of two among the ``\deleted{quarters}\inserted{slices}'' $Q_i$ ($i\in\{\unsure{C|0},1,\infty\}$) is a half-great circle passing in between $s_i$ and $s_j$ at equal
% archive foliation: 18
